\chapter{Input Models: State Preparation and Block Encoding}
% Solutions to the exercises of \cref{ch:lec06}.

\begin{solution}[Exercise~\ref{ex:lec06-1}]
Since $U_A$ is unitary, $\norm{U_A}=1$, and a compression of a unitary by isometries
cannot increase its norm:
\[
  \norm{A}=\norm{\alpha\,(\bra{0^a}\otimes\id)U_A(\ket{0^a}\otimes\id)}
  \le\alpha\,\norm{\bra{0^a}\otimes\id}\,\norm{U_A}\,\norm{\ket{0^a}\otimes\id}
  =\alpha\cdot1\cdot1\cdot1=\alpha,
\]
so $\alpha\ge\norm{A}$ is necessary. Conversely, if $\norm{A}\le\alpha$ then
$\norm{A/\alpha}\le1$, and \cref{prop:lec06-svd} gives a $(1,1)$-block-encoding
$U$ of $A/\alpha$, i.e.\ $(\bra{0}\otimes\id)U(\ket{0}\otimes\id)=A/\alpha$; reading
the same $U$ as an encoding of $A$ multiplies the block by $\alpha$, so $U$ is an
$(\alpha,1)$-block-encoding of $A$.

The max-norm bounds each entry but not the collective effect of a row: for
\[
  A=\begin{pmatrix}1&1\\1&-1\end{pmatrix},\qquad \norm{A}_{\max}=1,\quad
  \norm{A}=\sqrt2=s\,\norm{A}_{\max}\text{ with }s=2,
\]
because $A\dg A=2\id$ has both singular values equal to $\sqrt2$. This is the
$s$-sparse bound $\norm{A}\le s\norm{A}_{\max}$ saturated at $s=2$, and it explains
why the sparse construction pays $\alpha=s$: it must be safe for the worst row.
\end{solution}

\begin{solution}[Exercise~\ref{ex:lec06-2}]
By \cref{def:lec06-be} (exact case), $(\bra{0^a}\otimes\id)U_A(\ket{0^a}\otimes\id)
=A/\alpha$, so
\[
  U_A\ket{0^a}\ket{b}=\ket{0^a}\frac{A\ket{b}}{\alpha}+\ket{\perp},
  \qquad (\bra{0^a}\otimes\id)\ket{\perp}=0,
\]
where $\ket{\perp}$ lies entirely in the orthogonal complement of the flag
subspace. Measuring the flag register in the computational basis gives outcome
$0^a$ with probability
\[
  p=\norm{(\bra{0^a}\otimes\id)U_A\ket{0^a}\ket{b}}^2
   =\norm{\tfrac1\alpha A\ket{b}}^2=\frac{\norm{A\ket{b}}^2}{\alpha^2},
\]
and the post-measurement system state is the normalized flag component
$\bigl(A\ket{b}/\alpha\bigr)/\sqrt{p}=A\ket{b}/\norm{A\ket{b}}$. The expected number
of independent repetitions until the flag first reads $0^a$ is $1/p=
\alpha^2/\norm{A\ket{b}}^2$, i.e.\ it grows like $\alpha^2$. (With amplitude
amplification, \cref{sec:lec03-4}, this improves to $\bigo(\alpha/\norm{A\ket{b}})$
rounds.) An oversized $\alpha$---as in the $\alpha=10$ variant of
\cref{exa:lec06-2x2}---therefore inflates the runtime quadratically (or linearly,
after amplification) for no benefit, which is why $\alpha$ is kept as close to
$\norm{A}$ as the construction allows.
\end{solution}

\begin{solution}[Exercise~\ref{ex:lec06-3}]
The location map is $c(j)=(j+1)\bmod4$, a $4$-cycle, so
$O_c\ket{j}=\ket{(j+1)\bmod4}$ (equivalently the cyclic-shift permutation on two
qubits). The entry oracle rotates the signal qubit by the amplitude
$a_j=A_{c(j),j}$,
\[
  O_A\ket{0}\ket{i}\ket{j}
   =\bigl(a_j\ket{0}+\sqrt{1-a_j^2}\,\ket{1}\bigr)\ket{i}\ket{j},
\]
noting the amplitude depends on the column $j$ only (the nonzero row is determined by
$j$). With $U_A=(\id\otimes O_c)O_A$ (the signal qubit is untouched by $O_c$, which
acts on the column register, relabeling it as the row register),
\[
  U_A\ket{0}\ket{j}
  =\bigl(a_j\ket{0}+\sqrt{1-a_j^2}\ket{1}\bigr)\ket{(j+1)\bmod4},
\]
and therefore
\[
  \bra{0}\bra{i}U_A\ket{0}\ket{j}
  =a_j\,\braket{i|(j+1)\bmod4}=a_j\,\delta_{i,(j+1)\bmod4}=A_{ij},
\]
since $A_{ij}=a_j$ exactly when $i=(j+1)\bmod4$ and is $0$ otherwise. The $\ket{1}$
branch is annihilated by the $\bra{0}$ projection. By \cref{rem:lec06-verify} this
verifies $U_A\in\mathrm{BE}_{1,1}(A)$.
\end{solution}

\begin{solution}[Exercise~\ref{ex:lec06-4}]
For $A=\diag(a_0,a_1)$ the diagonal entry oracle
$O_A\ket{0}\ket{i}=(a_i\ket{0}+\sqrt{1-a_i^2}\ket{1})\ket{i}$ acts on the two-qubit
space (flag qubit tensor system qubit). Ordering the basis as
$\ket{0}\ket{0},\ket{0}\ket{1},\ket{1}\ket{0},\ket{1}\ket{1}$,
\[
  U_A=O_A=\begin{pmatrix}
    a_0 & 0 & \sqrt{1-a_0^2} & 0\\
    0 & a_1 & 0 & \sqrt{1-a_1^2}\\
    \sqrt{1-a_0^2} & 0 & -a_0 & 0\\
    0 & \sqrt{1-a_1^2} & 0 & -a_1
  \end{pmatrix},
\]
which \emph{is} symmetric and real here, hence Hermitian for this particular
convention of the controlled rotation. Whether the diagonal encoding comes out
Hermitian depends on the sign convention of the rotation; the substantive point is
that Hermiticity of $A$ places \emph{no} constraint on the off-diagonal blocks of
$U_A$ (the columns $\sqrt{1-a_i^2}$ and the sign of the lower-right block are free to
satisfy unitarity), so a generic block-encoding circuit for a Hermitian $A$---for
instance one that applies a phase $\ee^{\ii\phi_i}$ in the $\ket{1}$ branch, giving
\[
  \begin{pmatrix} a_i & \ee^{-\ii\phi_i}\sqrt{1-a_i^2}\\
                  \ee^{\ii\phi_i}\sqrt{1-a_i^2} & -a_i\ee^{\ii\phi_i}\cdots\end{pmatrix}
\]
in each block---breaks $U_A=U_A\dg$. Because qubitization (\cref{ch:lec08}) needs the
encoding unitary itself to be a reflection ($U_A^2=\id$), the Hermitian property must
be imposed as a separate design requirement, which is exactly \cref{def:lec06-hbe}.
\end{solution}

\begin{solution}[Exercise~\ref{ex:lec06-5}]
Trace the amplitudes in the proof of \cref{thm:lec06-ssparse}. The first diffusion
$D$ maps $\ket{0^{s'}}\mapsto\frac{1}{\sqrt s}\sum_\ell\ket{\ell}$, introducing a
factor $1/\sqrt s$ on the source side. On the target side, applying $D$ to
$\ket{0^{s'}}$ likewise produces $\frac{1}{\sqrt s}\sum_{\ell'}\ket{\ell'}$, a second
factor $1/\sqrt s$. When the inner product is formed, the two sums collapse to a
single surviving index via $\braket{\ell'|\ell}=\delta_{\ell\ell'}$, leaving the
product of the two prefactors, $\frac{1}{\sqrt s}\cdot\frac{1}{\sqrt s}=\frac1s$,
multiplying $A_{ij}$. Neither the entry oracle nor the location oracle contributes any
scalar (they are norm-preserving and act as permutations/controlled rotations), so the
$1/s$ comes \emph{entirely} from the two diffusion operators. Hence the top-left block
is $A/s$ and any such construction has $\alpha=s$.

For a dense matrix $s=N=2^n$, so $\alpha=N$, exponential in the qubit count.
Downstream algorithms cost $\bigo(\alpha\cdot\text{stuff})=\bigo(N)$ queries, wiping
out any quantum advantage: block encoding a generic dense matrix through the sparse
construction is no better than reading all $N^2$ entries. The construction is useful
precisely when $s=\polylog(N)$ (or the matrix is dense but structured, encodable by a
different route with small $\alpha$, as in the LCU or RACBEM models).
\end{solution}

\begin{solution}[Exercise~\ref{ex:lec06-6}]
Here $L=2$, $\beta_0=\beta_1=\tfrac12$, $\alpha=\beta_0+\beta_1=1$. The prepare
unitary must satisfy
$\mathrm{PREP}\ket0=\frac1{\sqrt\alpha}\sum_k\sqrt{\beta_k}\ket k
=\tfrac1{\sqrt2}(\ket0+\ket1)=\Hgate\ket0$, so $\mathrm{PREP}=\Hgate$, the single
flag-qubit Hadamard. The select unitary is
$\mathrm{SEL}=\proj0\otimes\Xgate+\proj1\otimes\Zgate$, applying $\Xgate$ when the
flag is $\ket0$ and $\Zgate$ when it is $\ket1$. Then, using $\Hgate\dg=\Hgate$,
\[
  (\bra0\otimes\id)U(\ket0\otimes\id)
  =(\bra0\Hgate\otimes\id)\,\mathrm{SEL}\,(\Hgate\ket0\otimes\id)
  =\sum_k\abs{\braket{k|\Hgate|0}}^2 V_k
  =\tfrac12\Xgate+\tfrac12\Zgate=A,
\]
since $\abs{\braket{k|\Hgate|0}}^2=\tfrac12$ for $k=0,1$. Thus
$U\in\mathrm{BE}_{1,1}(A)$. The spectral norm is
$\norm{A}=\tfrac12\norm{\Xgate+\Zgate}=\tfrac12\cdot\sqrt2=\tfrac1{\sqrt2}\approx0.707$
(the eigenvalues of $\Xgate+\Zgate$ are $\pm\sqrt2$). So $\alpha=1$ is \emph{not}
optimal: the LCU factor is the $\ell_1$-norm $\sum_k\beta_k=1$, which exceeds
$\norm{A}=1/\sqrt2$. LCU pays the $\ell_1$-norm, not the spectral norm; the two
coincide only when the $V_k$ act ``in the same direction,'' which $\Xgate$ and
$\Zgate$ (anticommuting) do not.
\end{solution}

\begin{solution}[Exercise~\ref{ex:lec06-7}]
Model on the Hermitian construction of \cref{thm:lec06-hbe}. Use two $n$-qubit
registers (row $i$, column $j$) and process the source with $V$ and the target with a
mirror copy of $V$. Applying $V$ to the column register of the source
$\ket{0^n}\ket{j}$ gives
\[
  V\ket{0^n}\ket{j}=\frac1{\sqrt N}\sum_i\ee^{\ii\theta_{ij}/2}\ket i\ket j .
\]
Applying $V$ to the target $\ket{0^n}\ket{i}$ gives
$\frac1{\sqrt N}\sum_{i'}\ee^{\ii\theta_{i'i}/2}\ket{i'}\ket i$. Inserting a
\textsc{swap} of the two registers between the source and target halves (so that the
row label produced on the source side is matched against the column label on the
target side) and taking the inner product,
\[
  \bra{0^n}\bra i\,U\,\ket{0^n}\ket j
  =\frac1N\sum_{k}\ee^{-\ii\theta_{ki}/2}\ee^{\ii\theta_{kj}/2}
   \braket{\text{matched}}
  =\frac1N\,\ee^{\ii(\theta_{ij})/2}\ee^{\ii\theta_{ij}/2}\cdot 1
  =\frac1N\ee^{\ii\theta_{ij}}=\frac{A_{ij}}{N},
\]
where Hermiticity $A_{ji}=A_{ij}^{\ast}$, i.e.\ $\theta_{ji}=-\theta_{ij}$, is used so
that the phase from the target copy of $V$ combines with the source phase to give
$\ee^{\ii\theta_{ij}}$ rather than canceling, and the $\braket{\text{matched}}$
Kronecker deltas from the \textsc{swap} pick out the single surviving term $k$ with
$k=i$ on one side and $k=j$ on the other. The two copies of $V$ each contribute a
$1/\sqrt N$ from the uniform superposition $V\ket{0^n}\ket\cdot$, so the top-left block
is $A/N$: the subnormalization factor is $\alpha=N$. (This matches the general
principle that an $N$-dense Hermitian matrix costs $\alpha=N$; the entries lying on
the unit circle is what makes the amplitudes $\ee^{\ii\theta_{ij}/2}/\sqrt N$ uniform
in magnitude and hence loadable by a single $V$.)
\end{solution}
